% !PRINTED 148
\weberpage{148}
%<<CONT
\[
 \nu_1,\nu_2,\nu_3\ldots\nu_n
\]
\gesperrt{form a decreasing, or at least never increasing, sequence of numbers}, or that the differences
\[
 \nu_1-\nu_2,\nu_2-\nu_3,\ldots\nu_{n-1}-\nu_n
\]
are all positive or at least not negative.

We conclude secondly that no term of the function $S$ can contain an exponent higher than $\nu_1$; for otherwise a term would also occur in which $\alpha_1$ had this higher exponent, and this would be of higher order than $A$, contrary to the supposition. Thus the terms that can occur at all in a symmetric function are, apart from their coefficients, completely determined by the highest term, and the number of possible terms is finite when the highest term is given.

We can use the exponent sequence of its highest term,
\[
 (\nu_1,\nu_2\ldots\nu_n)
\]
to denote the \gesperrt{order} of the symmetric function.

Apart from the case in which all exponents are zero, so that the function is a constant, the lowest possible order is
\[
 (1,0,0\ldots0),
\]
to which the single symmetric function
\[
 M(\alpha_1+\alpha_2+\cdots+\alpha_n)=-Ma_1
\]
corresponds.\begingroup\renewcommand{\thefootnote}{E2}\footnote[1]{[Source note:] These examples leave an arbitrary constant summand unstated; adding such a summand does not change the stated order. The formulas are retained as printed.}\endgroup{} The next higher order is
\[
 (1,1,0\ldots0),
\]
to which corresponds the symmetric function
\[
 Ma_1+M'a_2
\]
and thus in the first $n$ orders we obtain the fundamental functions themselves and their linear combinations, and no others.

The next order is characterized by
\[
 (2,0,0\ldots0)
\]
and contains the linear combinations of the fundamental functions with the sum of the squares.

In all these cases the representability of the symmetric functions has already been proved, and we shall therefore now give the general proof by deriving the representation
%CONT>>
